% G2-7 Kerangka pemikiran.
% Sumber: diolah penulis.
\begin{tikzpicture}[x=1mm,y=-1mm,font=\footnotesize,
    tahap/.style={draw,rectangle,text width=100mm,align=left,
      inner sep=3pt,minimum height=21mm},
    judul/.style={draw,rectangle,fill=black!10,text width=20mm,align=center,
      inner sep=3pt,minimum height=21mm,font=\small},
    panah/.style={-{Stealth[length=2.2mm]}}]

  % Tinggi kotak dan jarak antarbaris dibuat seragam: 21mm dan 27mm.
  \node[judul] (j1) at (12,0)   {Masalah};
  \node[tahap] (k1) at (76,0)
    {Harga BTC/IDR jangka pendek sulit diprediksi; pembanding Naive
     (harga tidak berubah) sulit dikalahkan.};

  \node[judul] (j2) at (12,27)  {Celah};
  \node[tahap] (k2) at (76,27)
    {Studi pada data menit umumnya memprediksi \textbf{arah}, bukan nilai
     harga; model boosting sering dipakai dengan setelan bawaan; jarang
     memakai pembanding naive dan uji Diebold-Mariano.};

  \node[judul] (j3) at (12,54)  {Pendekatan};
  \node[tahap] (k3) at (76,54)
    {Data candle 1 menit dan 5 menit BTC/IDR dari API Indodax; empat model
     boosting (XGBoost, LightGBM, CatBoost, AdaBoost.R2) $\times$ beberapa
     metode optimasi \textit{hyperparameter} $\times$ beberapa
     \textit{time frame} $\times$ tiga rasio \textit{split}
     (80:20, 70:30, 60:40).};

  \node[judul] (j4) at (12,81)  {Evaluasi};
  \node[tahap] (k4) at (76,81)
    {MAE, RMSE, MAPE, dan akurasi arah; uji Diebold-Mariano dengan koreksi
     HLN terhadap Naive.};

  \node[judul] (j5) at (12,108) {Keluaran};
  \node[tahap] (k5) at (76,108)
    {Jawaban atas RM1--RM3.};

  \foreach \a/\b in {1/2,2/3,3/4,4/5}{
    \draw[panah] (k\a.south) -- (k\b.north);
  }
\end{tikzpicture}
